\bookchapter{الفصل 3: اربط هاتفك بـ “خط الحياة” — الشبكة}

\chapterintro{بدون شبكة، يفقد الهاتف الذكي نصف ذكائه. الشبكة هي خط حياة الهاتف الحديث، فهمها واستخدامها جيداً وحمايتها هو فرض عين على كل مستخدم هاتف.}

\booksection{صفر: لا تضيع في الاتصال — ميّز “الطريق” و“السيارة” بين Wi-Fi والبيانات الخلوية}

الكثير من الناس يضيء أيقونتا Wi-Fi والإشارة في أعلى شاشة هاتفهم في وقت واحد، لكنهم في حيرة تامة من الفرق الجوهري بينهما. أكثر سوء فهم كلاسيكي هو:\textbf{الاتصال بـ Wi-Fi لكن لا يمكن تحميل الصفحات}— لأن الموجّه فقد الاتصال بالإنترنت، ولم يتحول الهاتف تلقائياً إلى البيانات الخلوية. لفهم هذه المشكلة بالكامل، نحتاج إلى بناء نموذج إدراكي “الطريق والسيارة”.

\begin{tcbinfo}{التشبيه الأساسي: الطريق والسيارة}
تخيل الإنترنت كمدينة.\textbf{“الطريق”}هو القناة بين هاتفك والإنترنت،\textbf{“السيارة”}هي البيانات التي تسير على الطريق. Wi-Fi والبيانات الخلوية طريقان مختلفان، لكن البيانات (السيارة) نفسها.

\end{tcbinfo}

\booksubsection{الفرق الجوهري: قناة محلية مقابل طريق سريع بعيد}

Wi-Fi في جوهرها\textbf{تقنية شبكة محلية}، فهو مثل الطريق المعبّد لاسلكياً تحت منزلك، ينقل البيانات من الموجّه إلى الهاتف. لكن إذا كانت فاتورة النطاق العريض غير مدفوعة أو انقطع الكابل البصري، فالطريق لا يزال موجوداً، لكن “الطريق السريع” الخارجي انقطع، وبالطبع لا يمكن الذهاب إلى أي مكان. أما البيانات الشبكية (البيانات المتنقلة) فتنقل عبر أبراج 4G/5G للمشغل، أي أنك تدخل مباشرة إلى الخط السريع للمشغل، طالما هناك إشارة يمكنك الاتصال المباشر بالإنترنت.

\booksubsection{التكلفة والسيناريو: اشتراك شهري ثابت مقابل دفع حسب الاستخدام}

Wi-Fi عادةً استثمار لمرة واحدة (رسوم النطاق العريض الشهرية)، طالما الكهرباء لا تنقطع، استخدام غير محدود للعائلة كلها، مناسب جداً لتنزيل الأفلام ولعب الألعاب وغيرها من السيناريوهات عالية الاستهلاك للبيانات. أما البيانات الخلوية فتُحسب حسب الاستخدام، وتجاوز الباقة قد يؤدي إلى خفض السرعة أو رسوم إضافية، أنسب للمهام الخفيفة في سيناريوهات التنقل مثل المترو وسفر العمل.

\begin{tcbpractice}{جرب: تحقق من باقة بياناتك}
افتح [الإعدادات] > [شبكة الهاتف المحمول] > [إدارة البيانات] في هاتفك، انظر كم البيانات المتبقية لك هذا الشهر، وأي التطبيقات تستهلك البيانات أكثر. كن على دراية، فلا تقلق عند الخروج.

\end{tcbpractice}

\booksubsection{الاعتبار الأمني: Wi-Fi العام “ميدان صيد” خفي}

من حيث الأمان، البيانات المتنقلة عادةً أكثر أماناً من Wi-Fi العام. البيانات المتنقلة مشفّرة من قبل المشغل، أما Wi-Fi العام فيمكن بسهولة أن ينصبه الهاكرز كـ “نقطة ساخنة مزيفة” لاعتراض البيانات غير المشفّرة. لكن يجب التمييز هنا: مع الانتشار الشامل لـ HTTPS اليوم، يمكن للمهاجم اعتراض\textbf{المواقع التي تزورها}(اسم النطاق)، و\textbf{محتوى HTTP غير المشفّر}؛ أما رسائل WeChat,كلمات مرور تسجيل الدخول لتطبيقات البنوك وغيرها من البيانات المشفّرة عبر HTTPS،\textbf{لا يمكن للمهاجم قراءتها مباشرة}— ما لم ينفذ هجوماً أكثر تعقيداً لفصل SSL وتجاهل المستخدم تحذيرات الشهادة. ومع ذلك، عند إجراء عمليات الدفع في الخارج، لا يزال استخدام بيانات الهاتف المحمول مباشرة هو الخيار الأكثر أماناً.

\textbf{مستوى متقدم: استخدم GitHub Pages لبناء مدونة شخصية مجانية}\par

في الأماكن العامة مثل المطارات والمولات والمقاهي، لا تتصل بشبكات Wi-Fi مجانية ذات أسماء مشبوهة (مثل “Free\_WiFi\_2024”). وحتى لو اتصلت، فلا تقم أبداً بتسجيل الدخول إلى تطبيقات البنوك أو إدخال كلمات المرور عبرها.

\booksubsection{الحكم من الهاتف: فهم “الرموز الخفية” في شريط الحالة}

كيف تعرف ما يستخدمه هاتفك الآن للاتصال بالإنترنت؟خذ هاتف Huawei كمثال: عند ظهور\textbf{أيقونة الأسهم ثنائية الاتجاه}بجانب إشارة المشغل أو بجانب إشارة WLAN في شريط الحالة، فهذا يعني أنك تستخدم الطريقة المقابلة للاتصال بالإنترنت؛ وإذا أضاءت كلتا الإشارتين معاً، فهذا يعني أن تسريع الشبكة قد دخل حيز التنفيذ، وأن الهاتف يستخدم القناتين معاً في الوقت نفسه.

أما مستخدمو iPhone فالأمر أكثر وضوحاً: إذا أضاءت أيقونة Wi-Fi فأنت تستخدم Wi-Fi، وإذا لم تُضئ فانظر إلى أعمدة الإشارة في الزاوية العلوية اليمنى لمعرفة أنك تستخدم بيانات الهاتف. كما تعرض بعض هواتف Android حروف “5G” أو “4G” في شريط الحالة، لتخبرك مباشرة بنوع شبكة الهاتف المحمول الحالية.

\booksection{أولاً: شبكة الهاتف المحمول (4G/5G) والبيانات}

شبكة الهاتف المحمول هي أساس اتصالك بالإنترنت في أي وقت ومكان. من 2G إلى 5G، ارتفعت السرعة آلاف المرات، لكن المبدأ يظل متشابهاً.

\booksubsection{معايير الشبكة:}

\begin{itemize}
  \item \textbf{4G LTE}: التغطية الأوسع حالياً، والسرعة كافية (نظرياً حتى 100Mbps+)، واستهلاك البطارية منخفض نسبياً.
  \item \textbf{5G}: سرعة عالية (نظرياً حتى 10Gbps)، لكن التغطية ما زالت قيد التحسين، واستهلاك البطارية مرتفع. مناسب لتنزيل الملفات الكبيرة ومشاهدة الفيديو عالي الدقة.
  \item \textbf{5G SA مقابل NSA}: SA هي شبكة قائمة بذاتها، أما NSA فتعتمد على محطات 4G. يوفر SA تجربة أفضل، لكنه يتطلب دعم المشغل والهاتف معاً.
\end{itemize}

\booksubsection{إدارة البيانات:}

\begin{itemize}
  \item \textbf{ما هي البيانات}: هي حجم البيانات الناتج عن استخدام الهاتف للإنترنت. مشاهدة الفيديو وتصفح الصور تستهلك أكبر قدر من البيانات، بينما النصوص والصوتيات تستهلك قدراً أقل.
  \item \textbf{نصائح لتوفير البيانات}: فعّل "وضع توفير البيانات" في الإعدادات، وقم بتحديث التطبيقات أو مشاهدة الفيديو عالي الدقة تحت Wi-Fi.
  \item \textbf{أعمدة الإشارة $\neq$ سرعة الشبكة}: قوة الإشارة تدل فقط على جودة الاتصال، أما سرعة الشبكة فتعتمد أيضاً على مدى انشغال المحطة وبُعدك عنها.
\end{itemize}

\textbf{تحذير مهم: أوقف[تجوال البيانات]قبل السفر للخارج!}\par

قد يؤدي استخدام بيانات الهاتف المحمول في الخارج إلى فواتير باهظة جداً! قبل السفر، تأكد من إيقاف تجوال البيانات في الإعدادات، أو اشترِ مسبقاً بطاقة بيانات محلية أو باقة تجوال دولية.

\booksection{ثانياً: Wi-Fi: الاتصال والمشاركة والأمان}

Wi-Fi هو الوسيلة الرئيسية للاتصال بالإنترنت في المنازل والمكاتب، سريع ولا يستهلك البيانات. لكن Wi-Fi يتطور باستمرار، وتختلف المعايير اختلافاً كبيراً.

\booksubsection{الاتصال بـ Wi-Fi:}

اذهب إلى [الإعدادات] > [WLAN] واختر الشبكة وأدخل كلمة المرور. تدعم العديد من الهواتف الآن الاتصال عبر مسح الرمز البرمجي، وهو أكثر راحة.

\booksubsection{مشاركة كلمة المرور:}

يمكن لكل من Android و iOS للهواتف المتصلة مشاركة كلمة المرور بسرعة مع الأصدقاء. ينشئ Android رمزاً برمجياً، بينما يمكن لـ iOS المشاركة تلقائياً عند التقارب.

\booksubsection{تطور معايير Wi-Fi:}

تُوضع معايير Wi-Fi من قبل IEEE، ولكل جيل تحسن ملحوظ:

\begin{center}
\begin{tabularx}{\linewidth}{|X|X|X|X|X|}
\hline
\textbf{المعيار} & \textbf{البروتوكول} & \textbf{النطاق الترددي} & \textbf{أقصى سرعة} & \textbf{الميزات} \tabularnewline
\hline
Wi-Fi 4 & 802.11n & 2.4/5GHz & 600Mbps & معيار قديم، التغطية مقبولة \tabularnewline
Wi-Fi 5 & 802.11ac & 5GHz & 3.5Gbps & الأكثر انتشاراً حالياً، السرعة كافية \tabularnewline
Wi-Fi 6 & 802.11ax & 2.4/5GHz & 9.6Gbps & أداء قوي مع الأجهزة المتعددة، زمن استجابة منخفض \tabularnewline
Wi-Fi 6E & 802.11ax & 2.4/5/6GHz & 9.6Gbps & إضافة نطاق 6GHz، أقل ازدحاماً \tabularnewline
Wi-Fi 7 & 802.11be & 2.4/5/6GHz & 46Gbps & نقل فائق السرعة، تجميع متعدد الروابط \tabularnewline
\hline
\end{tabularx}
\end{center}

\begin{itemize}
  \item \textbf{التقنيات الرئيسية في Wi-Fi 6}: يتيح OFDMA (تقسيم الترددات المتعامدة متعدد الوصول) لأجهزة متعددة النقل في الوقت نفسه، ويعزّز MU-MIMO (متعدد المستخدمين متعدد المداخل والمخارج) كفاءة التزامن بين الأجهزة المتعددة، ويسمح TWT (وقت الاستيقاف المستهدف) للأجهزة بتوفير الطاقة.
  \item \textbf{Wi-Fi 6E}: يضيف نطاق 6GHz إلى Wi-Fi 6، مما يوفر قنوات أكثر ويقلل التداخل. يتطلب دعم الهاتف والموجّه معاً.
  \item \textbf{Wi-Fi 7}: إدخال MLO (تشغيل متعدد الروابط)، يمكن استخدام عدة نطاقات لنقل البيانات في وقت واحد، مما يقلل التأخير بشكل كبير. قناة فائقة العرض 320MHz تمنح قفزة في السرعة. منصات الشرائح الرائدة مثل Snapdragon 8 Elite Gen5 وDimensity 9400 تدعم Wi-Fi 7، وهواتف مثل Xiaomi 17 وHuawei Mate 80 وSamsung S26 يمكنها استخدام شبكات Wi-Fi 7 (تتطلب موجّه Wi-Fi 7).
  \item \textbf{2.4GHz مقابل 5GHz}: 2.4GHz يخترق الجدران جيداً لكنه بطيء وبه تداخلات كثيرة؛ 5GHz سريع لكنه ضعيف في اختراق الجدران. نطاق 6GHz في Wi-Fi 6E/7 هو الأسرع لكنه الأضعف في اختراق الجدران، ويناسب الاستخدام داخل الغرفة نفسها.
\end{itemize}

\booksubsection{مخاطر Wi-Fi العام:}

قد يتعرض Wi-Fi العام للاختراق من قبل قراصنة يتنصتون على البيانات أو ينشئون نقاطاً ساخنة احتيالية. القواعد:

\begin{itemize}
  \item تجنب تسجيل الدخول إلى الحسابات البنكية أو إدخال كلمات المرور.
  \item الأفضل استخدام VPN لتشفير الاتصال.
  \item لا تتصل تلقائياً بشبكات Wi-Fi غير مألوفة.
  \item استخدم مواقع HTTPS (التي تظهر أيقونة القفل في شريط العنوان).
\end{itemize}

\booksection{ثالثاً: بلوتوث: اللاعب الشامل للاتصال اللاسلكي قصير المدى}

قد يكون بلوتوث أحد أكثر التقنيات اللاسلكية استخداماً على هاتفك — سماعات الرأس اللاسلكية، والساعات الذكية، وأنظمة السيارة، ولوحات المفاتيح والفأرة، كلها تعتمد عليه للاتصال.

\booksubsection{المبدأ التقني:}

يعمل بلوتوث في النطاق الترددي ISM 2.4GHz، ويستخدم تقنية الترددات القافزة نطاق الانتشار (FHSS) لتوزيع البيانات على 79 قناة والتبديل السريع بينها، مما يقي من التداخل فعّالاً. تنقسم مسافة اتصال بلوتوث إلى ثلاث فئات: Class 1 نحو 100 متر، وClass 2 نحو 10 أمتار (شائعة في الهواتف)، وClass 3 نحو متر واحد.

\booksubsection{تطور الإصدارات:}

\begin{itemize}
  \item \textbf{بلوتوث 4.0}: أدخل بلوتوث منخفض الطاقة (BLE)، مما أتاح للأجهزة ذات البطاريات الصغيرة مثل الأساور وأحزمة قياس نبضات القلب الاتصال اللاسلكي، وأطال عمر البطارية من أيام إلى أشهر.
  \item \textbf{بلوتوث 4.2}: حسّن سرعة النقل وحماية الخصوصية، ودعم الاتصال المباشر بـ IPv6.
  \item \textbf{بلوتوث 5.0}: ضاعف سرعة النقل (2Mbps)، وضاعف مسافة النقل أربع مرات (300 متر نظرياً)، وضاعف قدرة البث ثماني مرات. يدعم الصوت المزدوج (الاتصال بسماعتين في الوقت نفسه).
  \item \textbf{بلوتوث 5.1}: أدخل وظيفة توجيه الاتجاه (AoA/AoD)، بدقة تحديد موقع تصل إلى مستوى السنتيمتر.
  \item \textbf{بلوتوث 5.2}: معيار LE Audio الجديد، يدعم ترميز LC3 (جودة صوت أعلى واستهلاك طاقة أقل)، والصوت متعدد التدفقات، والصوت الإذاعي.
  \item \textbf{بلوتوث 5.3/5.4}: تحسينات إضافية في استهلاك الطاقة واستقرار الاتصال، وتعزيز البث الدوري.
\end{itemize}

\textbf{انتشار LE Audio}: في 2025-2026، دخل LE Audio مرحلة انتشار سريع. سماعات True Wireless الجديدة من العلامات الرئيسية (مثل Huawei FreeBuds Pro 5 وXiaomi Buds 6 Pro وSamsung Galaxy Buds3 Pro وغيرها) تدعم جميعها LE Audio وترميز LC3، مما يوفر تأخيراً أقل وجودة صوت أفضل. عند شراء سماعات جديدة، يُنصح بإعطاء الأولوية للمنتجات الداعمة لـ LE Audio.

\booksubsection{سيناريوهات التطبيق:}

\begin{itemize}
  \item \textbf{نقل الصوت}: سماعات الرأس اللاسلكية، وأنظمة صوت السيارة. أحدث LE Audio في بلوتوث 5.2 قفزة نوعية في جودة الصوت وزمن الاستجابة.
  \item \textbf{نقل البيانات}: نقل الملفات والصور عن قرب. رغم أن سرعته أقل من Wi-Fi Direct، إلا أنه يمتاز بالبساطة.
  \item \textbf{إقران الأجهزة}: الساعات الذكية، والأساور، ولوحات المفاتيح والفأرة، وأذرع التحكم بالألعاب.
  \item \textbf{خدمات الموقع}: منارات بلوتوث (Beacon) للملاحة الداخلية وتتبع الأشياء (مثل Apple AirTag).
  \item \textbf{المنزل الذكي}: شبكة بلوتوث Mesh للتحكم في المصابيح والأقفال وغيرها.
\end{itemize}

\booksection{رابعاً: التحكم بالأشعة تحت الحمراء: تحويل الهاتف إلى جهاز تحكم عن بُعد شامل}

التحكم بالأشعة تحت الحمراء من أقدم تقنيات التحكم اللاسلكي، لكنه ما زال عملياً على الهواتف. تضم بعض هواتف Android (مثل بعض موديلات Xiaomi وHuawei) مرسلاً بالأشعة تحت الحمراء مدمجاً، يمكنه التحكم في التلفزيونات والمكيفات ومستقبلات البث وغيرها من الأجهزة المنزلية.

\booksubsection{مبدأ العمل:}

يُصدر التحكم بالأشعة تحت الحمراء نبضات ضوئية بتردد محدد (عادة 38kHz) عبر صمام ثنائي باعث للأشعة تحت الحمراء، ويفكّك مستقبل الأشعة تحت الحمراء في الجهاز هذه النبضات لتنفيذ العملية المقابلة. تستخدم العلامات المختلفة بروتوكولات ترميز مختلفة، ومن الشائعة: NEC وRC5 وRC6 وSONY وغيرها. يحتوي تطبيق الأشعة تحت الحمراء في الهاتف على مكتبات ترميز لكل العلامات، ما عليك سوى اختيار العلامة والموديل للمطابقة.

\booksubsection{القيود:}

\begin{itemize}
  \item \textbf{الاتجاهية}: الأشعة تحت الحمراء تنتقل في خطوط مستقيمة، ويجب توجيهها نحو الجهاز، ولا تخترق الجدران.
  \item \textbf{المسافة القصيرة}: المسافة الفعّالة عادة 5-8 أمتار.
  \item \textbf{عدم دعم Apple}: لا يحتوي iPhone على مرسل بالأشعة تحت الحمراء، ولا يمكنه استخدام هذه الميزة.
\end{itemize}

\booksection{خامساً: NFC: سحر اللمس}

NFC (Near Field Communication، الاتصال قريب المدى) تقنية اتصال لاسلكي عالية التردد لمسافة قصيرة جداً (أقل من 10 سنتيمترات)، تعمل على تردد 13.56MHz، وبسرعة نقل تصل إلى 424Kbps.

\booksubsection{مبدأ الاتصال والأوضاع:}

تطوّرت NFC من تقنية RFID، وتدعم ثلاثة أوضاع تشغيل:

\begin{itemize}
  \item \textbf{وضع القراءة/الكتابة}: يقرأ الهاتف المعلومات من علامات NFC، مثل قراءة ملصق NFC لتشغيل إجراء.
  \item \textbf{وضع النظير إلى النظير}: جهازان NFC ينقلان البيانات بينهما، مثل Android Beam (متوقف).
  \item \textbf{وضع محاكاة البطاقة}: يحاكي الهاتف بطاقة ذكية، وهو الوضع الأكثر استخداماً — بطاقات المواصلات وبطاقات الدخول والبطاقات البنكية تستخدم هذا الوضع.
\end{itemize}

\booksubsection{آليات الأمان:}

\begin{itemize}
  \item \textbf{المسافة الفائقة القصر}: مسافة الاتصال ضمن 10 سنتيمترات هي بحد ذاتها ضمان أمان مادي، مما يجعل التنصت عن بُعد شبه مستحيل.
  \item \textbf{الشريحة الآمنة}: يحتوي الهاتف على شريحة أمان مستقلة (SE) أو يستخدم HCE (محاكاة البطاقة المضيفة)، وتُخزَّن مفاتيح الدفع في المنطقة الآمنة، ولا يمكن استخراجها حتى لو تم عمل root للهاتف.
  \item \textbf{التحقق بالبصمة/الوجه}: يتطلب الدفع تأكيداً بالتعرف البيومتري، فلا يمكن الاحتيال حتى لو فُقد الهاتف.
  \item \textbf{الترميز بالرموز}: لا تنقل Apple Pay وHuawei Pay وغيرها رقم البطاقة الحقيقي، بل تستخدم رمزاً مؤقتاً لا يمكن إعادة استخدامه حتى لو تم اعتراضه.
\end{itemize}

\booksubsection{التطبيقات النموذجية:}

\begin{itemize}
  \item \textbf{بطاقة المواصلات/بطاقة الدخول}: كتابة معلومات البطاقة المادية إلى الهاتف، ثم تمرير الهاتف يكفي.
  \item \textbf{الدفع عبر الهاتف}: Apple Pay وHuawei Pay وغيرها، تربط البطاقة البنكية بالهاتف، والدفع باللمس.
  \item \textbf{إقران الأجهزة}: لمس NFC لإقران سماعات بلوتوث ومكبرات الصوت بسرعة.
  \item \textbf{التحقق من الهوية}: قراءة بطاقة الهوية بـ NFC، والتحقق من رخصة القيادة الإلكترونية.
\end{itemize}

\booksection{سادساً: نقطة الاتصال الشخصية: تحويل الهاتف إلى موجّه Wi-Fi}

عندما يكون لديك بيانات وفيرة، يمكنك تشغيل نقطة اتصال للأجهزة اللوحية والكمبيوتر.

\begin{itemize}
  \item \textbf{طريقة التشغيل}: الإعدادات > نقطة الاتصال الشخصية/مشاركة الشبكة، وفعّلها.
  \item \textbf{تحذيرات}: تستهلك الكثير من البطارية والبيانات، تذكر تعيين كلمة مرور معقدة، وتحديد عدد الأجهزة المتصلة.
  \item \textbf{المشاركة عبر USB}: استخدام كابل بيانات للاتصال بالكمبيوتر، أكثر استقراراً وأسرع من نقطة الاتصال اللاسلكية.
\end{itemize}

\begin{tcbpractice}{جرّب بنفسك: إعداد نقطة اتصال وإدارتها}
\begin{enumerate}
  \item افتح [الإعدادات]، وابحث عن "نقطة الاتصال الشخصية" أو "مشاركة الشبكة"، وفعّلها، وعيّن الاسم وكلمة المرور.
  \item استخدم جهازاً آخر (مثل الكمبيوتر) للبحث عن شبكة Wi-Fi هذه والاتصال بها.
  \item تحقق من الهاتف ما إذا كان يمكنك إدارة الأجهزة المتصلة، أو طرد الأجهزة غير المعروفة.
\end{enumerate}

\end{tcbpractice}
